\thispagestyle{empty}
\begin{center}
\vspace*{38mm}
{\displayfont\Huge Épitaphes chrétiennes\\[2mm] datées de Rome\\[4mm]}
{\displayfont\Large (340–403)\par}
\vspace{10mm}
{\large\itshape Transmission numérique, chronologie et formulaire\\ d’un corpus consulaire\par}
\vfill
\end{center}
\cleardoublepage

\thispagestyle{empty}
\begin{center}
\vspace*{20mm}
{\large\scshape Alexis Angel Martini\par}
\vspace{22mm}
{\displayfont\Huge Épitaphes chrétiennes\\[2mm] datées de Rome\\[4mm]}
{\displayfont\Large (340–403)\par}
\vspace{10mm}
{\large\itshape Transmission numérique, chronologie et formulaire\\ d’un corpus consulaire\par}
\vspace{18mm}
{\color{prismered}\rule{40mm}{0.4pt}}\par
\vspace{6mm}
{\normalsize Catalogue raisonné de \st{nA} notices,\\ études critiques, concordances et index\par}
\vfill
{\small Manuscrit soumis à évaluation scientifique\\ Version 4.0 — 14 octobre 2026\par}
\vspace{6mm}
{\small\scshape Projet PRISME\par}
\end{center}
\clearpage

\thispagestyle{empty}
\noindent{\small
\textbf{Alexis Angel Martini}, \emph{Épitaphes chrétiennes datées de Rome (340–403). Transmission numérique, chronologie et formulaire d’un corpus consulaire}. Manuscrit, version 4.0, 14 octobre 2026.

\medskip
\noindent\textbf{Statut du document.} Manuscrit complet préparé pour une évaluation scientifique par les pairs. Il ne constitue ni une publication, ni une édition critique des inscriptions, ni un bon à tirer. Son acceptation dépend des rapporteurs et de l’éditeur.

\medskip
\noindent\textbf{Citation.} Auteur, titre, version 4.0, numéro de notice (n\textsuperscript{o}) ou page. Les numéros de notice 1–\st{nA} sont stables d’une version à l’autre.

\medskip
\noindent\textbf{Données et programmes.} Les données consolidées (\texttt{corpus\_v4.json}), les scripts de calcul et de composition, le classeur de collation EDCS/EDB et le registre des consultations accompagnent le manuscrit (dépôt de l’auteur, dossier \texttt{07\_monographie\_v4}). Toutes les valeurs chiffrées du texte sont produites par ces scripts ; voir l’annexe.

\medskip
\noindent\textbf{Données de tiers.} Les extraits et données descriptives de l’Epigraphic Database Bari (EDB) sont reproduits avec attribution, selon les licences indiquées par la base (fichiers EpiDoc : CC BY-NC-SA 3.0 ; pages de consultation : CC BY-SA 4.0 ; notice du dépôt Zenodo : CC BY 4.0). Les textes de l’Epigraphik-Datenbank Clauss/Slaby (EDCS) sont cités d’après les exports diffusés par le projet SDAM (université d’Aarhus). Aucune photographie n’est reproduite.

\medskip
\noindent\textbf{Déclaration sur l’usage d’outils d’intelligence artificielle.} La préparation de ce manuscrit a fait appel à des modèles de langage, sous la direction et la responsabilité de l’auteur : extraction et structuration des données, écriture des programmes de contrôle et de calcul, collation automatisée entre bases, aide à la rédaction et composition typographique. Les contrôles effectués, leurs limites et les sources effectivement consultées sont décrits dans le volume. Cette déclaration doit être adaptée à la politique de l’éditeur auquel le manuscrit est soumis.

\medskip
\noindent\textbf{Composition.} XeLaTeX ; caractères Libertinus (SIL Open Font License) et FreeSerif pour quelques signes diacritiques.
}
\cleardoublepage
